\documentclass[a4paper,11pt]{article}

\usepackage[margin=25mm]{geometry}
\usepackage{amsmath}
\usepackage{amssymb}
\usepackage{array}
\usepackage{booktabs}

\title{Starting Methods and Starting Torque of Induction Motors}
\author{}
\date{}

\begin{document}

\maketitle

\section{Basic Concept}

The starting method of an induction motor depends on the type of rotor
and the type of power supply.

At the instant of starting, the rotor speed is zero:

\[
N = 0.
\]

Therefore, the slip is

\[
s = \frac{N_s-N}{N_s} = 1,
\]

where \(N_s\) is the synchronous speed and \(N\) is the rotor speed.

Thus,

\[
\boxed{s=1 \quad \text{at starting}}
\]

is an important relation.

\section{Three-Phase Wound-Rotor Induction Motor}

A wound-rotor induction motor has three-phase windings on its rotor.
These rotor windings are connected to external resistors through
\textbf{slip rings} and brushes.

\[
\boxed{
\text{Rotor winding}
\rightarrow
\text{Slip rings}
\rightarrow
\text{Brushes}
\rightarrow
\text{External resistance}
}
\]

At starting,

\[
s=1.
\]

By increasing the rotor resistance \(R_2\), the slip at which the
maximum torque occurs moves toward \(s=1\).

In simplified form,

\[
s_{\mathrm{max}} \propto R_2.
\]

Therefore,

\[
R_2 \uparrow
\quad\Rightarrow\quad
s_{\mathrm{max}} \rightarrow 1
\quad\Rightarrow\quad
T_{\mathrm{start}} \uparrow.
\]

Hence, the purpose of adding external rotor resistance is to obtain
a large starting torque.

\[
\boxed{
\text{Wound rotor}
\rightarrow
\text{Increase rotor resistance}
\rightarrow
\text{Large starting torque}
}
\]

The slip rings themselves do not increase the torque.
Their purpose is to provide an electrical connection between the
rotating rotor winding and the stationary external resistance.

\section{Three-Phase Squirrel-Cage Induction Motor}

In a squirrel-cage induction motor, the rotor circuit cannot be
connected to an external resistance.

Therefore, starting is controlled from the \textbf{stator (power-supply) side}.

Typical starting methods are:

\begin{itemize}
    \item Star-delta starting
    \item Autotransformer starting
    \item Inverter starting
\end{itemize}

For an induction motor, the starting torque is approximately
proportional to the square of the applied voltage:

\[
\boxed{T_{\mathrm{start}} \propto V^2}.
\]

Therefore, if the starting voltage is reduced,

\[
V \downarrow
\quad\Rightarrow\quad
T_{\mathrm{start}} \downarrow.
\]

However, the starting current is also reduced.

\subsection{Star-Delta Starting}

The motor is connected in star (\(Y\)) during starting and changed to
delta (\(\Delta\)) during normal operation.

For the same line voltage \(V_L\), the phase voltage during star
connection is

\[
V_{\mathrm{ph},Y}
=
\frac{V_L}{\sqrt{3}}.
\]

Since

\[
T \propto V^2,
\]

the starting torque becomes approximately

\[
T_Y
=
\left(
\frac{1}{\sqrt{3}}
\right)^2 T_{\Delta}
=
\frac{1}{3}T_{\Delta}.
\]

Thus,

\[
\boxed{
T_{Y,\mathrm{start}}
\approx
\frac{1}{3}
T_{\Delta,\mathrm{start}}
}
\]

for the same line voltage.

\subsection{Autotransformer Starting}

An autotransformer is inserted between the power supply and the motor
during starting.

It reduces the motor terminal voltage:

\[
V_{\mathrm{motor}} < V_{\mathrm{supply}}.
\]

Therefore,

\[
V \downarrow
\quad\Rightarrow\quad
I_{\mathrm{start}} \downarrow,
\qquad
T_{\mathrm{start}} \downarrow.
\]

\subsection{Inverter Starting}

An inverter controls both the voltage and frequency supplied to the
motor.

\[
\boxed{
\text{Inverter}
\rightarrow
V \text{ control} + f \text{ control}
}
\]

By starting at a low frequency and appropriately controlling the
voltage, the motor can start smoothly while producing useful torque.

\section{Single-Phase Induction Motor}

A single-phase induction motor cannot produce a starting torque using
only its main winding.

At standstill,

\[
\boxed{T_{\mathrm{start}}=0}
\]

for an ideal single-phase induction motor with only the main winding.

Therefore, an auxiliary winding is used.

The current in the auxiliary winding is made to have a different
\textbf{phase} from the current in the main winding.

\[
I_{\mathrm{main}}
\quad\text{and}\quad
I_{\mathrm{aux}}
\]

have a phase difference:

\[
\phi \neq 0.
\]

This phase difference produces a rotating magnetic field component,
which creates starting torque.

Thus,

\[
\boxed{
\text{Single phase}
\rightarrow
\text{Auxiliary winding}
\rightarrow
\text{Phase difference}
\rightarrow
\text{Starting torque}
}
\]

\section{Summary}

\begin{center}
\renewcommand{\arraystretch}{1.5}
\begin{tabular}{p{0.25\textwidth} p{0.32\textwidth} p{0.32\textwidth}}
\toprule
\textbf{Motor Type} &
\textbf{Starting Method} &
\textbf{Starting-Torque Concept} \\
\midrule

Three-phase wound rotor &
External rotor resistance through slip rings &
Increase \(R_2\) so that the maximum-torque point approaches
\(s=1\), producing a large starting torque. \\

Three-phase squirrel cage &
Star-delta, autotransformer, or inverter &
The rotor circuit cannot be externally modified.
For voltage-reduction starting,
\(T_{\mathrm{start}}\propto V^2\). \\

Single-phase induction motor &
Auxiliary winding with a phase difference &
Main winding alone gives zero starting torque ideally.
A phase difference produces starting torque. \\

\bottomrule
\end{tabular}
\end{center}

\section{Key Relations to Remember}

\[
\boxed{s=1 \qquad \text{at starting}}
\]

\[
\boxed{
\text{Wound rotor: }
R_2\uparrow
\Rightarrow
s_{\mathrm{max}}\rightarrow 1
\Rightarrow
T_{\mathrm{start}}\uparrow
}
\]

\[
\boxed{
\text{Squirrel cage: }
T_{\mathrm{start}}\propto V^2
}
\]

\[
\boxed{
\text{Single phase: }
T_{\mathrm{start}}=0
\quad
\text{with the main winding alone}
}
\]

\[
\boxed{
\text{Auxiliary winding}
+
\text{phase difference}
\Rightarrow
\text{starting torque}
}
\]

\end{document}
